\documentclass[preview]{standalone}

\usepackage{amsmath}
\usepackage{amssymb}
\usepackage{stellar}
\usepackage{definitions}
\usepackage{bettelini}

\begin{document}

\id{probability-continuous-distributions}
\genpage

\section{Uniform distribution}

\begin{snippetdefinition}{continuous-uniform-distribution-definition}{Uniform distribution on an interval}
    Let \(a,b\in\realnumbers\) with \(a<b\). A real-valued
    \randomvariable \(X\) has the \emph{uniform distribution} on \([a,b]\),
    written
    \[
        X\sim\operatorname{Unif}([a,b]),
    \]
    if \(X\) is an \absolutelycontinuousrandomvariable with density
    \[
        f_X(x)=\frac{1}{b-a}\indicator_{[a,b]}(x).
    \]
\end{snippetdefinition}

\begin{snippetproposition}{continuous-uniform-distribution-cdf}{Distribution function of a uniform random variable}
    Let \(a,b\in\realnumbers\) with \(a<b\), and let
    \(X\sim\operatorname{Unif}([a,b])\). Then
    \[
        F_X(x)
        =
        \begin{cases}
            0 & x<a,\\
            \dfrac{x-a}{b-a} & a\leq x\leq b,\\
            1 & x>b.
        \end{cases}
    \]
\end{snippetproposition}

\begin{snippetproof}{continuous-uniform-distribution-cdf-proof}{continuous-uniform-distribution-cdf}{Distribution function of a uniform random variable}
    Integrate the density:
    \[
        F_X(x)=\int_{-\infty}^x\frac{1}{b-a}\indicator_{[a,b]}(t)\,dt.
    \]
    The three cases correspond to \(x<a\), \(a\leq x\leq b\), and \(x>b\).
\end{snippetproof}

\begin{snippetproposition}{continuous-uniform-distribution-mean-variance}{Mean and variance of a uniform random variable}
    Let \(a,b\in\realnumbers\) with \(a<b\), and let
    \(X\sim\operatorname{Unif}([a,b])\). Then
    \[
        \expectation[X]=\frac{a+b}{2},
        \qquad
        \variance(X)=\frac{(b-a)^2}{12}.
    \]
\end{snippetproposition}

\begin{snippetproof}{continuous-uniform-distribution-mean-variance-proof}{continuous-uniform-distribution-mean-variance}{Mean and variance of a uniform random variable}
    The expected value is
    \[
        \expectation[X]
        =
        \frac{1}{b-a}\int_a^b x\,dx
        =
        \frac{a+b}{2}.
    \]
    Also,
    \[
        \expectation[X^2]
        =
        \frac{1}{b-a}\int_a^b x^2\,dx
        =
        \frac{a^2+ab+b^2}{3}.
    \]
    Therefore
    \[
        \variance(X)
        =
        \expectation[X^2]-\expectation[X]^2
        =
        \frac{a^2+ab+b^2}{3}-\frac{(a+b)^2}{4}
        =
        \frac{(b-a)^2}{12}.
    \]
\end{snippetproof}

\section{Exponential distribution}

\begin{snippetdefinition}{exponential-distribution-definition}{Exponential distribution}
    Let \(\lambda>0\). A real-valued \randomvariable \(X\) has the
    \emph{exponential distribution} with rate \(\lambda\), written
    \[
        X\sim\operatorname{Exp}(\lambda),
    \]
    if \(X\) is an \absolutelycontinuousrandomvariable with density
    \[
        f_X(x)=\lambda e^{-\lambda x}\indicator_{[0,+\infty)}(x).
    \]
\end{snippetdefinition}

\begin{snippetproposition}{exponential-distribution-cdf-tail}{Distribution function and tail of an exponential random variable}
    Let \(\lambda>0\), and let \(X\sim\operatorname{Exp}(\lambda)\). Then
    \[
        F_X(x)
        =
        \begin{cases}
            0 & x<0,\\
            1-e^{-\lambda x} & x\geq0,
        \end{cases}
    \]
    and, for every \(x\geq0\),
    \[
        \mathbb P(X>x)=e^{-\lambda x}.
    \]
\end{snippetproposition}

\begin{snippetproof}{exponential-distribution-cdf-tail-proof}{exponential-distribution-cdf-tail}{Distribution function and tail of an exponential random variable}
    If \(x<0\), then the density vanishes on \((-\infty,x]\). If
    \(x\geq0\), then
    \[
        F_X(x)
        =
        \int_0^x\lambda e^{-\lambda t}\,dt
        =
        1-e^{-\lambda x}.
    \]
    Hence
    \[
        \mathbb P(X>x)=1-F_X(x)=e^{-\lambda x}.
    \]
\end{snippetproof}

\begin{snippettheorem}{exponential-distribution-memoryless-property}{Memoryless property}
    Let \(\lambda>0\), and let \(X\sim\operatorname{Exp}(\lambda)\). For
    every \(y,h\geq0\),
    \[
        \mathbb P(X>y+h\mid X>y)
        =
        \mathbb P(X>h).
    \]
\end{snippettheorem}

\begin{snippetproof}{exponential-distribution-memoryless-property-proof}{exponential-distribution-memoryless-property}{Memoryless property}
    Since \(\mathbb P(X>y)>0\),
    \[
        \mathbb P(X>y+h\mid X>y)
        =
        \frac{\mathbb P(X>y+h)}{\mathbb P(X>y)}
        =
        \frac{e^{-\lambda(y+h)}}{e^{-\lambda y}}
        =
        e^{-\lambda h}
        =
        \mathbb P(X>h).
    \]
\end{snippetproof}

\begin{snippetproposition}{exponential-distribution-mean-variance}{Mean and variance of an exponential random variable}
    Let \(\lambda>0\), and let \(X\sim\operatorname{Exp}(\lambda)\). Then
    \[
        \expectation[X]=\frac{1}{\lambda},
        \qquad
        \variance(X)=\frac{1}{\lambda^2}.
    \]
\end{snippetproposition}

\begin{snippetproof}{exponential-distribution-mean-variance-proof}{exponential-distribution-mean-variance}{Mean and variance of an exponential random variable}
    Integration by parts gives
    \[
        \expectation[X]
        =
        \int_0^{+\infty}\lambda xe^{-\lambda x}\,dx
        =
        \frac{1}{\lambda}.
    \]
    A second integration by parts gives
    \[
        \expectation[X^2]
        =
        \int_0^{+\infty}\lambda x^2e^{-\lambda x}\,dx
        =
        \frac{2}{\lambda^2}.
    \]
    Hence
    \[
        \variance(X)
        =
        \expectation[X^2]-\expectation[X]^2
        =
        \frac{2}{\lambda^2}-\frac{1}{\lambda^2}
        =
        \frac{1}{\lambda^2}.
    \]
\end{snippetproof}

\end{document}
